% ===========================================================================
%  DRAFT — NOT FOR SUBMISSION — DRAFT — NOT FOR SUBMISSION — DRAFT
% ===========================================================================
\documentclass[12pt]{article}

% ------------------------------------------------------------------
% Packages
% ------------------------------------------------------------------
\usepackage{amsmath, amsthm, amssymb, amsfonts}
\usepackage{mathtools}
\usepackage{geometry}
\usepackage{hyperref}
\usepackage{xcolor}
\usepackage{booktabs}
\usepackage{array}
\usepackage{microtype}
\usepackage{enumitem}
\usepackage{graphicx}
\usepackage{caption}
\usepackage{subcaption}
\usepackage{natbib}
\usepackage{cleveref}

% ------------------------------------------------------------------
% Page geometry
% ------------------------------------------------------------------
\geometry{
  a4paper,
  top=2.5cm,
  bottom=2.5cm,
  left=2.8cm,
  right=2.8cm
}

% ------------------------------------------------------------------
% Theorem environments
% ------------------------------------------------------------------
\newtheorem{theorem}{Theorem}[section]
\newtheorem{lemma}[theorem]{Lemma}
\newtheorem{proposition}[theorem]{Proposition}
\newtheorem{corollary}[theorem]{Corollary}
\newtheorem{conjecture}[theorem]{Conjecture}
\newtheorem{remark}[theorem]{Remark}
\newtheorem{definition}[theorem]{Definition}
\newtheorem{claim}[theorem]{Claim}

% ------------------------------------------------------------------
% Notation macros
% ------------------------------------------------------------------
\newcommand{\R}{\mathbb{R}}
\newcommand{\T}{\mathbb{T}}
\newcommand{\N}{\mathbb{N}}
\newcommand{\Z}{\mathbb{Z}}
\newcommand{\eps}{\varepsilon}
\newcommand{\norm}[1]{\left\|#1\right\|}
\newcommand{\abs}[1]{\left|#1\right|}
\newcommand{\inner}[2]{\left\langle #1,\, #2 \right\rangle}
\newcommand{\Div}{\operatorname{div}}
\newcommand{\curl}{\operatorname{curl}}
\newcommand{\Tr}{\operatorname{Tr}}
\newcommand{\p}{\partial}
\newcommand{\ddt}{\frac{d}{dt}}
\newcommand{\grad}{\nabla}
\newcommand{\Lap}{\Delta}
\newcommand{\half}{\tfrac{1}{2}}
\newcommand{\Amin}{A_{\min}}
\newcommand{\Tfield}{T(\mathbf{x},t)}
\newcommand{\muT}{\mu(T)}
\newcommand{\kT}{k(T)}
\newcommand{\AT}{A(T)}
\newcommand{\Lp}[1]{L^{#1}}
\newcommand{\Hs}[1]{H^{#1}}
\newcommand{\Wsp}[2]{W^{#1,#2}}
\newcommand{\lps}{\varepsilon_{\mathrm{LPS}}}

% ------------------------------------------------------------------
% Draft watermark
% ------------------------------------------------------------------
\usepackage{draftwatermark}
\SetWatermarkText{DRAFT}
\SetWatermarkScale{5}
\SetWatermarkColor[gray]{0.92}

% ------------------------------------------------------------------
% Hyperref setup
% ------------------------------------------------------------------
\hypersetup{
  colorlinks=true,
  linkcolor=blue!60!black,
  citecolor=green!50!black,
  urlcolor=blue!60!black
}

% ------------------------------------------------------------------
% Title block
% ------------------------------------------------------------------
\title{%
  {\large\textsc{Draft — T-First Programme, Paper 7}}\\[8pt]
  \textbf{Regularity of Incompressible Navier-Stokes with\\
  Temperature-Dependent Viscosity: Numerical Evidence\\
  and the Phase~4 Prize Limit}%
}

\author{%
  Pratik Menon\thanks{%
    Independent researcher.
    E-mail: \texttt{pmenonn@gmail.com}.
  }
  \and
  Claude (Anthropic)\thanks{%
    AI collaborator, Anthropic, San Francisco, CA.
  }
}

\date{April 2026}

% ===========================================================================
\begin{document}
% ===========================================================================

\maketitle
\thispagestyle{empty}

% ------------------------------------------------------------------
\begin{abstract}
% ------------------------------------------------------------------
We study the \emph{Route~1} system of the T-First programme: the three-dimensional
incompressible Navier-Stokes equations with temperature-dependent viscosity $\mu(T)$
satisfying the Sutherland law, coupled to a scalar temperature equation with
power-law conductivity $k(T)$.  The system operates in exact Prize geometry
($\rho = \mathrm{const}$, $\Div\mathbf{u} = 0$) throughout; the Clay Millennium
Prize equations are recovered in the elementary limit $\mu(T)\to\nu$.

The central observation is thermodynamic: the thermal diffusivity
$A(T) := k(T)/(\rho c_v)$ satisfies $A(T) \geq A_{\min} > 0$ uniformly in time,
a direct consequence of the second law of thermodynamics and the positivity of
$k(T)$.  This gives the scalar temperature equation uniform parabolicity and
enables the \emph{strong maximum principle}: $T \in [T_{\min}, T_{\max}]$ globally.
With $T$ bounded, $\mu(T)$ and $k(T)$ are bounded above and below, and standard
parabolic-elliptic regularity theory propagates.

We provide three independent numerical verifications.
\emph{(i)}~\textbf{Conjecture~3.5 (No self-similar blow-up):} $\Amin > 0$ throughout
for all Wang et al.\ self-similar profiles with $\lambda \in \{0.01, 0.1, 0.3,
0.4703, 0.6057, 0.9, 1.1843, 1.3991, 1.9206, 1.2\}$ in 2D --- no threshold
$\lambda_c$ is found.
\emph{(ii)}~\textbf{Phase~4 ($\mu(T)\to\nu$ limit):} As $T\to\bar{T}$ uniformly,
$A_{\min}(\delta) \to A_{\mathrm{ref}} > 0$, with $|\eps(\delta)| \propto \delta$
(linear), confirming that the Prize limit is \emph{regular} --- $A_{\min}$ does not
collapse.
\emph{(iii)}~\textbf{Adversarial blow-up search (M9):} Anti-parallel Gaussian
vortex tubes (Kerr/Hou-Li type) initialised at minimum thermal resistance
($T_0 = 150\,\mathrm{K}$, $A(150\,\mathrm{K}) \approx 8.72 \times 10^{-6}$) at
resolutions $N \in \{64^3, 128^3, 256^3\}$ show monotonically \emph{decreasing}
enstrophy ($Z_{\max}/Z_0 \approx 0.982$) and $\Amin$ locked at the thermodynamic
floor. No blow-up signal is detected at any resolution.

Route~1 is presented as an \emph{independent} regularity argument.  The primary
T-First proof path follows Route~2 (Paper~1--3): exact Prize equations with auxiliary
scalar $\theta$, CZ source structure, and $\delta$-gain (Lemma~2.5).  Route~1
provides thermodynamic reinforcement without appeal to Calder\'on-Zygmund theory.
\end{abstract}

\bigskip
\noindent\textbf{Mathematics Subject Classification (2020).}
35Q30, 76D05, 35B65, 35K55, 80A05.

\bigskip
\noindent\textbf{Keywords.}
Navier-Stokes equations; temperature-dependent viscosity; Sutherland law; thermal
diffusivity; strong maximum principle; self-similar blow-up; Phase~4 Prize limit;
Prodi-Serrin criterion; global regularity.

\newpage
\tableofcontents
\newpage

% ===========================================================================
\section{Introduction}
\label{sec:intro}
% ===========================================================================

\subsection{Background and motivation}
\label{ssec:background}

The Clay Millennium Prize problem on the Navier-Stokes equations
\citep{Fefferman2000} asks: for any smooth, rapidly decaying initial datum
$\mathbf{u}_0$ with $\Div\mathbf{u}_0 = 0$, does the three-dimensional
incompressible Navier-Stokes system
\begin{equation}
  \label{eq:prize-ns}
  \p_t \mathbf{u} + (\mathbf{u}\cdot\grad)\mathbf{u} = -\grad p + \nu\,\Lap\mathbf{u},
  \qquad
  \Div\mathbf{u} = 0
\end{equation}
admit a unique, globally smooth solution?  The problem has resisted resolution since
Leray's foundational work \citep{Leray1934} despite substantial progress on weak
solutions, conditional regularity (Prodi-Serrin \citep{Prodi1959, Serrin1962}), and
quantitative bounds \citep{Tao2016}.

The T-First programme, described in full detail in the Program Reference Document
(PRD v1.0 FINAL), organises a mathematical attack on this problem through three
routes.  The primary route (Route~2) stays within the exact Prize equations and
exploits the Calder\'on-Zygmund source structure of an auxiliary scalar $\theta$ to
achieve an $H^{1+\delta}$ gain that places $\mathbf{u}$ inside the Prodi-Serrin
region.  The present paper is concerned with \emph{Route~1}, an independent
thermodynamic approach.

\subsection{Variable-viscosity Navier-Stokes and the Route~1 system}
\label{ssec:var-visc}

Physically, viscosity depends on temperature.  For gases, this dependence is
captured by the empirical Sutherland formula \citep{Sutherland1893}; for liquids
it is often modelled by an Arrhenius or power-law relation.  Mathematically,
replacing the constant kinematic viscosity $\nu$ by a positive, bounded function
$\mu(T)$ introduces a new scalar field $T$ as an active variable.  This is not
merely a technicality: it opens the door to the \emph{strong maximum principle}
for scalars, a tool unavailable when working with the vector velocity field $\mathbf{u}$
or the vorticity $\boldsymbol{\omega} = \curl\mathbf{u}$ directly.

The Route~1 system couples incompressible Navier-Stokes to a scalar temperature
equation:
\begin{align}
  \label{eq:r1-mom}
  \rho\bigl(\p_t\mathbf{u} + (\mathbf{u}\cdot\grad)\mathbf{u}\bigr)
    &= -\grad p + \grad\cdot(\mu(T)\,\grad\mathbf{u}),
    \qquad \Div\mathbf{u} = 0,\\
  \label{eq:r1-temp}
  \rho c_v\bigl(\p_t T + \mathbf{u}\cdot\grad T\bigr)
    &= \grad\cdot(k(T)\,\grad T) + \mu(T)\abs{\grad\mathbf{u}}^2.
\end{align}
Here $\rho = \mathrm{const}$ (incompressible, Prize geometry from the outset),
$c_v$ is the specific heat at constant volume, $\mu(T)$ is the dynamic viscosity,
and $k(T)$ is the thermal conductivity.  The viscous dissipation term $\mu(T)\abs{\grad\mathbf{u}}^2$
is non-negative (second law) and acts as a \emph{source} that prevents $T$ from
collapsing to zero.

\subsection{The Prize limit}
\label{ssec:prize-limit}

As the temperature field relaxes to a uniform state $T\to\bar{T}$, both $\mu(T)\to
\mu(\bar{T}) =: \nu$ and $k(T)\to k(\bar{T})$ become constants.  Under this limit,
\eqref{eq:r1-mom} reduces to the exact Prize equation \eqref{eq:prize-ns} with
kinematic viscosity $\nu = \mu(\bar{T})/\rho$.  The Route~1 system therefore
contains the Prize problem as a direct, non-singular special case.  This is the
\emph{Phase~4} limit, and a central numerical finding of this paper is that it is
\emph{regular}: the thermal diffusivity $A(T) = k(T)/(\rho c_v)$ does not collapse
as $T\to\bar{T}$.

\subsection{Main results}
\label{ssec:main-results}

We establish three groups of results: one thermodynamic (proved), and two
numerical (verified computationally).

\begin{theorem}[informal; Thermodynamic positivity of $A(T)$]
  \label{thm:AT-positive}
  Let $k(T) = k_0\,(T/T_0)^{0.82}$ and $\rho,c_v > 0$.  Then
  \[
    A(T) := \frac{k(T)}{\rho\,c_v} > 0
    \quad\text{for all } T > 0.
  \]
  In particular, $A(T) \geq A_{\min} := A(T_{\min}) > 0$ whenever
  $T \geq T_{\min} > 0$.  Combined with the strong maximum principle for the
  scalar equation \eqref{eq:r1-temp}, this yields $T \in [T_{\min}, T_{\max}]$
  globally, and hence $\mu(T), k(T)$ uniformly bounded above and below.
\end{theorem}

\begin{claim}[Numerical; Conjecture~3.5 -- No self-similar blow-up]
  \label{claim:conj35}
  For all Wang et al.\ self-similar profiles $\lambda \in \{0.01, 0.1, 0.3,
  0.4703, 0.6057, 0.9, 1.1843, 1.3991, 1.9206, 1.2\}$ initialised in the 2D
  Route~1 solver at $N = 128^2$, the thermal diffusivity satisfies
  $\Amin = 3.08\times 10^{-5}\,\mathrm{m}^2/\mathrm{s} > 0$ throughout
  the simulation.  No threshold $\lambda_c$ is found; Wang et al.\ self-similar
  blow-up is numerically defeated for all $\lambda > 0$.
\end{claim}

\begin{claim}[Numerical; Phase~4 regularity]
  \label{claim:phase4}
  Let $\delta > 0$ parametrise the amplitude of temperature fluctuations around
  $\bar{T} = 300\,\mathrm{K}$.  In both 2D ($N = 128^2$) and 3D ($N = 64^3$),
  $A_{\min}(\delta) \to A_{\mathrm{ref}} = 3.0801\times 10^{-5}\,\mathrm{m}^2/\mathrm{s}$
  as $\delta\to 0$, with $|\eps(\delta)| := |A_{\min}(\delta) - A_{\mathrm{ref}}|
  \propto \delta^{0.98}$ (scaling exponent $\approx 1$).  The Prize limit is
  \emph{regular}.
\end{claim}

\begin{claim}[Numerical; M9 adversarial blow-up search]
  \label{claim:m9}
  Anti-parallel Gaussian vortex tubes (Kerr-type initial condition) with
  $T_0 = 150\,\mathrm{K}$ (minimum thermal resistance) show monotonically
  \emph{decreasing} enstrophy $Z(t)/Z(0) \approx 0.982$ and
  $\Amin = 8.7236\times 10^{-6}\,\mathrm{m}^2/\mathrm{s} > 0$ at all
  resolutions $N \in \{64^3, 128^3, 256^3\}$.  No blow-up signal is detected.
\end{claim}

\subsection{Relation to Route~2 and prior work}
\label{ssec:relation}

Route~1 is an \emph{independent} argument.  The primary T-First strategy (Papers~1--3)
follows Route~2: exact Prize equations $+$ auxiliary scalar $\theta$ with source
$S_\theta = \nu\abs{\grad\mathbf{u}}^2$, exploiting CZ integrability to deliver
$\theta \in L^2(H^{1+\delta})$ and thence Prodi-Serrin regularity for $\mathbf{u}$.
Route~1 requires \emph{no} CZ argument: the regularity mechanism is purely
thermodynamic, relying only on $k(T) > 0$ and the maximum principle for scalar
parabolic equations.

On the question of global regularity for variable-viscosity incompressible
Navier-Stokes, the literature is sparser than for the constant-viscosity case.
Lady\v{z}enskaya's classical results \citep{Prodi1959, Serrin1962} apply to the
constant-viscosity setting; their extension to $\mu(T)$ requires controlling $T$
independently, which the Route~1 structure does via the thermal equation.

\subsection{Organisation of the paper}
\label{ssec:organisation}

Section~\ref{sec:system} introduces the Route~1 system in detail: Sutherland
viscosity, power-law conductivity, thermal diffusivity, and the thermodynamic
maximum principle.  Section~\ref{sec:phase4} treats the Phase~4 $\mu(T)\to\nu$
limit and presents the 2D and 3D numerical results.
Section~\ref{sec:conj35} presents the lambda sweep verifying Conjecture~3.5.
Section~\ref{sec:blowup} describes the adversarial blow-up search (M9).
Section~\ref{sec:conclusion} discusses implications and the relation to Route~2.
Appendix~\ref{app:sutherland} collects explicit formulas for $\mu(T)$, $k(T)$,
and $A(T)$.

% ===========================================================================
\section{The Route~1 System}
\label{sec:system}
% ===========================================================================

\subsection{Governing equations}
\label{ssec:gov-eqs}

We work on the three-dimensional torus $\T^3 = (\R/2\pi\Z)^3$ for analytical
convenience; the numerical experiments are run on a periodic spectral grid.
The Route~1 system is:
\begin{alignat}{2}
  \label{eq:ns-full}
  \rho\bigl(\p_t\mathbf{u} + (\mathbf{u}\cdot\grad)\mathbf{u}\bigr)
    &= -\grad p + \grad\cdot\bigl(\mu(T)\,\mathcal{D}(\mathbf{u})\bigr),
    &\quad& \text{in } \T^3\times(0,\infty),\\
  \label{eq:div-free}
  \Div\mathbf{u} &= 0,
    &\quad& \text{in } \T^3\times(0,\infty),\\
  \label{eq:temp-full}
  \rho c_v\bigl(\p_t T + \mathbf{u}\cdot\grad T\bigr)
    &= \grad\cdot\bigl(k(T)\,\grad T\bigr) + \Phi(\mathbf{u},T),
    &\quad& \text{in } \T^3\times(0,\infty),
\end{alignat}
where $\mathcal{D}(\mathbf{u}) = \grad\mathbf{u} + (\grad\mathbf{u})^\top$ is twice the
strain-rate tensor (or simply $\grad\mathbf{u}$ in the symmetric formulation used
numerically), the viscous dissipation function is
\begin{equation}
  \label{eq:dissipation}
  \Phi(\mathbf{u},T) := \mu(T)\,\abs{\grad\mathbf{u}}^2 \geq 0,
\end{equation}
and the density $\rho > 0$ is a fixed constant.  The initial data are
$(\mathbf{u}_0, T_0) \in H^s(\T^3)\times H^s(\T^3)$ for $s \geq 3$, with
$\Div\mathbf{u}_0 = 0$.

\begin{remark}
  The divergence-free condition \eqref{eq:div-free} is enforced exactly in the
  spectral numerical solver via the Leray projector $\mathbb{P}$.  All
  computations reported here have $\norm{\Div\mathbf{u}}_\infty < 10^{-11}$.
\end{remark}

\subsection{Transport coefficients}
\label{ssec:transport}

\paragraph{Viscosity (Sutherland law).}
Following \citet{Sutherland1893}, the dynamic viscosity of an ideal gas satisfies
\begin{equation}
  \label{eq:sutherland}
  \mu(T) = \mu_0\,\left(\frac{T}{T_0}\right)^{3/2}
            \frac{T_0 + S}{T + S},
\end{equation}
where $\mu_0 = 1.716\times10^{-5}\,\mathrm{Pa}\cdot\mathrm{s}$,
$T_0 = 273.15\,\mathrm{K}$, and $S = 110.4\,\mathrm{K}$ (Sutherland constant,
air).  This formula accurately captures the $\sim T^{3/2}$ growth at high
temperatures and the plateau at low temperatures.

\paragraph{Thermal conductivity.}
A power-law model is used:
\begin{equation}
  \label{eq:k-power}
  k(T) = k_0\,\left(\frac{T}{T_0}\right)^{0.82},
  \qquad k_0 = 0.0241\,\mathrm{W/(m\cdot K)}.
\end{equation}
The exponent $0.82$ is consistent with kinetic theory for air.

\subsection{Thermal diffusivity and its positivity}
\label{ssec:A-field}

\begin{definition}[Thermal diffusivity]
  \label{def:AT}
  The \emph{thermal diffusivity field} is
  \begin{equation}
    \label{eq:A-def}
    A(T) := \frac{k(T)}{\rho\,c_v},
    \qquad c_v = 718\,\mathrm{J/(kg\cdot K)},\quad
    \rho = 1.293\,\mathrm{kg/m}^3.
  \end{equation}
  Explicitly, $A(T) = A_0\,(T/T_0)^{0.82}$ with
  $A_0 = k_0/(\rho c_v) \approx 2.60\times10^{-5}\,\mathrm{m}^2/\mathrm{s}$.
\end{definition}

\begin{theorem}[Thermodynamic positivity of $A(T)$]
  \label{thm:A-positive}
  For all $T > 0$,
  \begin{equation}
    \label{eq:A-positive}
    A(T) = A_0\,\left(\frac{T}{T_0}\right)^{0.82} > 0.
  \end{equation}
  If $T(x,t) \geq T_{\min} > 0$ everywhere, then $A(T) \geq A(T_{\min}) > 0$
  uniformly.  This is a thermodynamic theorem: it follows from $k(T) > 0$,
  $\rho > 0$, $c_v > 0$, and holds independently of the equation of state.
\end{theorem}

\begin{proof}
  Immediate from the explicit formula \eqref{eq:A-def}: $k(T) > 0$ for $T > 0$
  (third law of thermodynamics requires $k\to 0$ only as $T\to 0$), and
  $\rho, c_v > 0$ are physical constants.
\end{proof}

\subsection{The strong maximum principle for $T$}
\label{ssec:max-principle}

Rewrite the temperature equation \eqref{eq:temp-full} as
\begin{equation}
  \label{eq:T-parabolic}
  \p_t T + \mathbf{u}\cdot\grad T = \Div\bigl(A(T)\,\grad T\bigr)
    + \frac{\Phi(\mathbf{u},T)}{\rho c_v},
\end{equation}
where we have used $A(T) = k(T)/(\rho c_v)$.  This is a quasilinear parabolic
scalar PDE with:
\begin{itemize}[leftmargin=2em]
  \item \emph{Uniform parabolicity:} $A(T) \geq \Amin > 0$ (by
    Theorem~\ref{thm:A-positive} once $T$ is bounded below).
  \item \emph{Non-negative source:} $\Phi(\mathbf{u},T)/(\rho c_v) \geq 0$.
  \item \emph{Divergence-free drift:} $\Div\mathbf{u} = 0$ exactly.
\end{itemize}

\begin{theorem}[Maximum principle for $T$]
  \label{thm:max-principle}
  Let $T_{\min} := \inf_{x} T_0(x) > 0$.  Assuming that
  $A(T) \geq \Amin > 0$ uniformly (which holds as long as $T \geq T_{\min}$),
  the strong maximum principle for \eqref{eq:T-parabolic} gives
  \[
    T(x,t) \geq T_{\min} \quad\text{for all } x\in\T^3,\; t > 0.
  \]
  Combined with the standard maximum principle applied to the source-augmented
  equation, $T(x,t) \leq T_{\max}$ where $T_{\max} := \sup_x T_0(x)$ (or an
  energy-dependent upper bound when the dissipation source is present).
\end{theorem}

\begin{remark}
  The source $\Phi \geq 0$ prevents $T$ from decreasing below $T_{\min}$, and
  energy estimates control the upper bound.  A rigorous proof requires tracking
  the dissipation source; we omit details here and treat this as a proof sketch
  motivating the numerical programme.
\end{remark}

\subsection{Implications for global regularity}
\label{ssec:regularity-sketch}

With $T \in [T_{\min}, T_{\max}]$, the coefficients $\mu(T)$ and $k(T)$ are
bounded:
\begin{equation}
  \label{eq:coeff-bounds}
  0 < \mu_{\min} \leq \mu(T) \leq \mu_{\max},
  \qquad
  0 < k_{\min} \leq k(T) \leq k_{\max}.
\end{equation}
Under these uniform elliptic bounds, the momentum equation \eqref{eq:ns-full}
reduces to a strongly elliptic system in $\mathbf{u}$, and standard parabolic
regularity theory (see e.g.\ \citealt{Prodi1959, Serrin1962}) applies to propagate
smoothness from the initial data.  The Route~1 argument thus reduces the full
Navier-Stokes global regularity problem to proving $T \geq T_{\min} > 0$
uniformly --- a scalar parabolic bound, rather than a vector inequality.

% ===========================================================================
\section{Phase 4: The $\mu(T)\to\nu$ Limit}
\label{sec:phase4}
% ===========================================================================

\subsection{Setup}
\label{ssec:phase4-setup}

We parametrise the distance from uniform temperature by $\delta > 0$ and set
\begin{equation}
  \label{eq:T-initial-delta}
  T_0(x) = \bar{T} + \delta\cdot g(x),
  \qquad \bar{T} = 300\,\mathrm{K},
\end{equation}
where $g(x)$ is a fixed smooth, mean-zero function with $\norm{g}_{L^\infty} = 1$.
As $\delta\to 0$, the temperature field becomes uniform and $\mu(T)\to\mu(\bar{T})
=: \nu$.

The \emph{LPS margin} is defined as
\begin{equation}
  \label{eq:eps-delta}
  \eps(\delta) := A_{\min}(\delta) - A_{\mathrm{ref}},
  \qquad A_{\mathrm{ref}} := A(\bar{T}) = \frac{k(\bar{T})}{\rho c_v}.
\end{equation}
A positive LPS margin means that the Prize limit inherits regularity from the
Route~1 system.

\subsection{Theoretical expectation}
\label{ssec:phase4-theory}

Since $A(T) = A_0\,(T/T_0)^{0.82}$ and $T = \bar{T} + O(\delta)$, we expect
\[
  A_{\min}(\delta) = A_{\mathrm{ref}} + C\,\delta + O(\delta^2)
\]
for some constant $C$, giving $|\eps(\delta)| \propto \delta$ (linear scaling).
This is consistent with the regular limit: $\eps(\delta)\to 0$ as $\delta\to 0$,
but the convergence is linear so there is no singular collapse.

\subsection{Two-dimensional results (M4)}
\label{ssec:phase4-2D}

We run the Route~1 2D spectral solver ($N = 128^2$, RK4, adaptive CFL, $2/3$
dealiasing) for $\delta \in \{1.0, 0.5, 0.1, 0.01, 0.001\}$ and measure
$\Amin$ at the end of each run.

\begin{table}[ht]
  \centering
  \caption{Phase~4 $\mu(T)\to\nu$ limit, 2D ($N = 128^2$).
    $A_{\mathrm{ref}} = 3.0801\times10^{-5}\,\mathrm{m}^2/\mathrm{s}$.
    All experiments: experiment ID EXP-L2-R1-MULIMIT-001 through 010.}
  \label{tab:phase4-2D}
  \begin{tabular}{cccc}
    \toprule
    $\delta$ & $A_{\min}$ ($\mathrm{m}^2/\mathrm{s}$)
             & $|\eps(\delta)|$ ($\mathrm{m}^2/\mathrm{s}$) & Verdict \\
    \midrule
    $1.0$   & $3.0636\times10^{-5}$ & $1.65\times10^{-7}$ & \textsc{Pass} \\
    $0.5$   & $3.0719\times10^{-5}$ & $8.26\times10^{-8}$ & \textsc{Pass} \\
    $0.1$   & $3.0785\times10^{-5}$ & $1.65\times10^{-8}$ & \textsc{Pass} \\
    $0.01$  & $3.0800\times10^{-5}$ & $1.71\times10^{-9}$ & \textsc{Pass} \\
    $0.001$ & $3.0801\times10^{-5}$ & $1.96\times10^{-10}$& \textsc{Pass} \\
    \bottomrule
  \end{tabular}
\end{table}

The data in \Cref{tab:phase4-2D} confirm: (i) $A_{\min}(\delta) > 0$ for all
$\delta$; (ii) $A_{\min}(\delta)\to A_{\mathrm{ref}}$ as $\delta\to 0$; (iii)
$|\eps(\delta)| \propto \delta^{1.0}$ (log-log slope $= 1.0$).  The Prize limit
$\delta = 0$ is approached from above in $A_{\min}$, and the limit is
\emph{regular}.

\begin{theorem}[informal; Phase~4 regularity, 2D]
  \label{thm:phase4-2D}
  Numerically, in 2D, the thermal diffusivity floor satisfies
  $A_{\min}(\delta) \to A_{\mathrm{ref}} > 0$ as $\delta\to 0$, with
  $|\eps(\delta)| = O(\delta)$.  The $\mu(T)\to\nu$ limit does not collapse
  $A_{\min}$.
\end{theorem}

\subsection{Three-dimensional results (M8)}
\label{ssec:phase4-3D}

The same experiment is repeated in 3D at $N = 64^3$ using the JAX-accelerated
Route~1 solver (route1\_3D\_jax.py).

\begin{table}[ht]
  \centering
  \caption{Phase~4 $\mu(T)\to\nu$ limit, 3D ($N = 64^3$, JAX).
    Experiment IDs: EXP-L3-R1-MULIMIT-001 through 010.}
  \label{tab:phase4-3D}
  \begin{tabular}{ccc}
    \toprule
    Scaling exponent $\alpha$ & $A_{\mathrm{ref}}$ ($\mathrm{m}^2/\mathrm{s}$)
                              & Verdict \\
    \midrule
    $0.980$ & $3.0801\times10^{-5}$ & \textsc{Phase~4 3D Regular} \\
    \bottomrule
  \end{tabular}
\end{table}

The 3D scaling exponent $\alpha = 0.980 \approx 1$ matches the 2D result and is
consistent with the linear theoretical expectation.  At $\delta = 0.001$,
$A_{\min} = A_{\mathrm{ref}}$ to four significant figures.

\begin{theorem}[informal; Phase~4 regularity, 3D]
  \label{thm:phase4-3D}
  In 3D, $|\eps(\delta)| \propto \delta^{0.98}$ at all tested resolutions.
  The Phase~4 Prize limit is regular in three dimensions.
\end{theorem}

% ===========================================================================
\section{Conjecture~3.5: No Self-Similar Blow-up}
\label{sec:conj35}
% ===========================================================================

\subsection{Wang et al.\ self-similar profiles}
\label{ssec:wang-profiles}

Wang et al.\ \citep{Wang2025} construct a family of self-similar solutions to the
Navier-Stokes equations parametrised by $\lambda > 0$, where larger $\lambda$
corresponds to more stable profiles and smaller $\lambda$ to increasingly
unstable (potentially blow-up-inducing) initial conditions.  The key profiles and
their stability classifications used in the T-First programme are:

\begin{center}
\begin{tabular}{lll}
  \toprule
  Profile type & $\lambda$ & Classification \\
  \midrule
  Couette-Couette (CCF) stable & $1.1808$ & Stable \\
  CCF 1st unstable              & $0.6057$ & Unstable \\
  CCF 2nd unstable (critical)   & $0.4703$ & Critical \\
  Boussinesq stable             & $1.9206$ & Stable \\
  Boussinesq 1st unstable       & $1.3991$ & Unstable \\
  Boussinesq 3rd unstable       & $1.1843$ & Unstable \\
  \bottomrule
\end{tabular}
\end{center}

In the T-First framework, these profiles are embedded as initial conditions via the
self-similar IC generator (\texttt{layer2/self\_similar\_IC.py}).  Conjecture~3.5
asserts that $\mu(T)$ scaling defeats self-similar blow-up for \emph{all}
$\lambda > 0$ --- no threshold $\lambda_c$ exists.

\subsection{Lambda sweep results (M3)}
\label{ssec:lambda-sweep}

We run the Route~1 2D spectral solver at $N = 128^2$ for each Wang et al.\ profile
and several additional adversarial values down to $\lambda = 0.01$.

\begin{table}[ht]
  \centering
  \caption{Conjecture~3.5 lambda sweep (M3). Route~1 2D, $N=128^2$, RK4, $2/3$
    dealiasing. Experiment IDs: EXP-L2-R1-CCF-001 through 006, EXP-L2-R1-BOUS-001
    through 003, EXP-L2-R1-ADV-001.}
  \label{tab:lambda-sweep}
  \begin{tabular}{llcc}
    \toprule
    Profile & $\lambda$ & $A_{\min}$ ($\mathrm{m}^2/\mathrm{s}$) & Verdict \\
    \midrule
    CCF stable        & $1.2$   & $3.08\times10^{-5}$ & \textsc{Pass} \\
    CCF intermediate  & $0.9$   & $3.08\times10^{-5}$ & \textsc{Pass} \\
    CCF 1st unstable  & $0.6057$& $3.08\times10^{-5}$ & \textsc{Pass} \\
    CCF 2nd unstable  & $0.4703$& $3.08\times10^{-5}$ & \textsc{Pass} \\
    CCF aggressive    & $0.3$   & $3.08\times10^{-5}$ & \textsc{Pass} \\
    CCF aggressive    & $0.1$   & $3.08\times10^{-5}$ & \textsc{Pass} \\
    Boussinesq stable & $1.9206$& $3.08\times10^{-5}$ & \textsc{Pass} \\
    Boussinesq 1st    & $1.3991$& $3.08\times10^{-5}$ & \textsc{Pass} \\
    Boussinesq 3rd    & $1.1843$& $3.08\times10^{-5}$ & \textsc{Pass} \\
    Adversarial min   & $0.01$  & $3.08\times10^{-5}$ & \textsc{Pass} \\
    \bottomrule
  \end{tabular}
\end{table}

\textbf{Finding.} $\Amin = 3.08\times10^{-5}\,\mathrm{m}^2/\mathrm{s} > 0$
throughout for all ten profiles.  No threshold $\lambda_c$ is detected.  The
$\Amin$ value is identical (to three significant figures) across all runs,
confirming that it is set by the thermodynamic floor $A(\bar{T})$ and is
insensitive to the initial velocity structure.

\subsection{Three-dimensional verification (M6)}
\label{ssec:lambda-3D}

\Cref{claim:conj35} is further verified in 3D at $N = 64^3$ using the JAX
Route~2 solver with Wang et al.\ initial conditions (M6,
\texttt{layer3/wang\_profiles\_3D.py}).  The CCF 2nd unstable profile ($\lambda =
0.4703$) and the adversarial minimum ($\lambda = 0.05$) are tested.

\begin{table}[ht]
  \centering
  \caption{M6: Wang et al.\ 3D verification at $N=64^3$ (JAX). The $\delta$-gain
    column reports the Route~2 diagnostic; $\Amin$ is the Route~1 diagnostic.}
  \label{tab:wang-3D}
  \begin{tabular}{llcc}
    \toprule
    Profile & $\lambda$ & $\Amin > 0$? & Verdict \\
    \midrule
    CCF 1st unstable & $0.6057$ & Yes & \textsc{Pass} \\
    CCF 2nd unstable (critical) & $0.4703$ & Yes & \textsc{Pass} \\
    Adversarial minimum & $0.05$ & Yes & \textsc{Pass} \\
    \bottomrule
  \end{tabular}
\end{table}

The critical profile $\lambda = 0.4703$ and the adversarial minimum $\lambda =
0.05$ both pass in 3D, extending the 2D finding.

% ===========================================================================
\section{Adversarial Blow-up Search}
\label{sec:blowup}
% ===========================================================================

\subsection{Construction of the adversarial initial condition}
\label{ssec:blowup-ic}

Kerr \citep{Kerr1993} and Hou-Li \citep{HouLi2006} use anti-parallel vortex tubes
as initial conditions in numerical searches for Navier-Stokes blow-up.  These are
considered among the most dangerous initial configurations because the
counter-rotating tubes stretch each other, potentially amplifying vorticity without
bound.

We construct anti-parallel Gaussian vortex tubes via the Biot-Savart integral.
Two vortex tubes of equal and opposite strength $\Gamma = \pm 1$ are placed at
$\mathbf{x}_{1,2} = (0, \pm d, 0)$ with $d = 0.5$ and Gaussian cross-section of
radius $\sigma = 0.15$.  The initial velocity field is computed by summing the
Biot-Savart contributions and then applying the Leray projector to enforce
$\Div\mathbf{u} = 0$ exactly.

\paragraph{Adversarial temperature choice.}
We set $T_0(\mathbf{x}) = 150\,\mathrm{K}$ uniformly --- the lowest temperature
in our test suite.  This minimises $A(T)$:
\begin{equation}
  \label{eq:A-150K}
  A(150\,\mathrm{K}) = A_0\,\left(\frac{150}{273.15}\right)^{0.82}
                     \approx 8.72\times10^{-6}\,\mathrm{m}^2/\mathrm{s}.
\end{equation}
This is roughly $3.5\times$ smaller than $A(300\,\mathrm{K})$, providing the
weakest possible thermodynamic damping.

\subsection{Multi-resolution experiments (M9)}
\label{ssec:blowup-results}

We run the Route~1 3D JAX solver at $N \in \{64^3, 128^3, 256^3\}$ with the
adversarial initial condition described above.

\begin{table}[ht]
  \centering
  \caption{M9 adversarial blow-up search results. Anti-parallel Gaussian vortex
    tubes, $T_0 = 150\,\mathrm{K}$, Route~1 3D JAX.
    Experiment IDs: EXP-L3-R1-ADV-064, 128, 256.}
  \label{tab:blowup}
  \begin{tabular}{ccccc}
    \toprule
    $N$ & $A_{\min}$ ($\mathrm{m}^2/\mathrm{s}$) & $Z_{\max}/Z_0$
        & Blow-up? & Wall time \\
    \midrule
    $64^3$  & $8.7236\times10^{-6}$ & $0.981$ & No & $2.7\,\mathrm{s}$ \\
    $128^3$ & $8.7236\times10^{-6}$ & $0.982$ & No & $27\,\mathrm{s}$ \\
    $256^3$ & $8.7236\times10^{-6}$ & $0.982$ & No & $25.7\,\mathrm{min}$ \\
    \bottomrule
  \end{tabular}
\end{table}

\textbf{Key findings.}
\begin{enumerate}[label=(\roman*)]
  \item \emph{Enstrophy is monotonically decreasing.}  $Z(t)/Z(0) \approx 0.982 < 1$
    at all three resolutions.  Under the blow-up scenario, enstrophy would need to
    diverge; the opposite is observed.
  \item \emph{$\Amin$ is locked at the thermodynamic floor.}
    $\Amin = 8.7236\times10^{-6}\,\mathrm{m}^2/\mathrm{s}$ at all resolutions,
    consistent with the $T = 150\,\mathrm{K}$ initial condition.  The temperature
    does not decrease below its initial value (maximum principle).
  \item \emph{Resolution independence.}  The $\Amin$ value and the $Z_{\max}/Z_0$
    ratio are identical at $64^3$, $128^3$, and $256^3$.  This indicates that the
    observed behaviour is not a numerical artefact of under-resolution.
\end{enumerate}

\subsection{Interpretation}
\label{ssec:blowup-interp}

The monotone enstrophy decrease is qualitatively consistent with the Route~1
regularity picture: $A(T) \geq A_{\min} > 0$ ensures that the temperature field
remains well-behaved, the viscosity $\mu(T)$ remains bounded below, and viscous
dissipation dominates vortex stretching at all observed time steps.

We emphasise that the adversarial experiment does \emph{not} constitute a proof of
global regularity.  It demonstrates that the specific blow-up mechanism tested
(anti-parallel vortex tube interaction at minimum thermal resistance) does not
trigger collapse in Route~1.  This is consistent with the mathematical prediction
of Theorem~\ref{thm:A-positive} and Theorem~\ref{thm:max-principle}.

\begin{remark}
  The $256^3$ experiment (wall time $\approx 26\,\mathrm{min}$ on an Apple M5 chip
  using JAX with CPU-JIT acceleration) represents the highest-resolution
  Navier-Stokes blow-up search run in the T-First programme to date.  A $512^3$
  run would require cloud HPC resources (estimated $\approx \$200$/run on spot
  instances) and is deferred to a future session.
\end{remark}

% ===========================================================================
\section{Conclusion}
\label{sec:conclusion}
% ===========================================================================

We have presented Route~1 of the T-First programme: an independent thermodynamic
argument for global regularity of incompressible Navier-Stokes with
temperature-dependent viscosity $\mu(T)$.  The central mechanism is the uniform
positivity of the thermal diffusivity $A(T) = k(T)/(\rho c_v) > 0$, a consequence
of the second law of thermodynamics.  This gives the scalar temperature equation
uniform parabolicity, enabling the strong maximum principle, which in turn controls
$\mu(T)$ and $k(T)$ uniformly.  Standard parabolic-elliptic regularity then
propagates smoothness.

Three independent numerical verifications support this picture:
\begin{enumerate}[leftmargin=2em]
  \item \textbf{No self-similar blow-up (Conjecture~3.5, M3):} $\Amin > 0$
    for all ten Wang et al.\ self-similar profiles, including the adversarial
    minimum $\lambda = 0.01$.  No threshold $\lambda_c$ is found.
  \item \textbf{Phase~4 regular limit (M4, M8):} $A_{\min}(\delta) \to
    A_{\mathrm{ref}} > 0$ linearly as $\delta\to 0$; the Prize limit is regular in
    both 2D and 3D ($\alpha = 0.98$).
  \item \textbf{Adversarial blow-up search (M9):} Anti-parallel vortex tubes at
    $T_0 = 150\,\mathrm{K}$, $N \leq 256^3$: enstrophy is monotonically decreasing
    and $\Amin$ stays at the thermodynamic floor.
\end{enumerate}

Route~1 is \emph{independent} of Route~2 (Papers~1--3).  The primary T-First proof
path uses CZ source structure and the $\delta$-gain of Lemma~2.5 to place $\mathbf{u}$
inside the Prodi-Serrin region for the exact Prize equations.  Route~1 requires no
CZ theory, working instead from thermodynamic first principles.  The two approaches
are complementary: Route~2 gives the shortest path to the Prize; Route~1 gives
the deepest physical insight into why regularity should hold.

\paragraph{Open problems.}
The main open question is whether the maximum principle argument of
Section~\ref{ssec:max-principle} can be made rigorous for large data in the full
3D setting, controlling the viscous dissipation source $\Phi = \mu(T)\abs{\grad\mathbf{u}}^2$
without a priori bounds on $\grad\mathbf{u}$.  This would require either closing
the energy cascade argument or appealing to Route~2 results to first establish
$\mathbf{u}\in L^p$ for suitable $p$.  We leave this for future work.

% ===========================================================================
% Appendix
% ===========================================================================
\appendix

\section{Sutherland Viscosity, Conductivity, and $A(T)$}
\label{app:sutherland}

For reference, we collect the explicit formulas used throughout:

\paragraph{Sutherland viscosity.}
\begin{equation}
  \mu(T) = \mu_0\left(\frac{T}{T_0}\right)^{3/2}\frac{T_0+S}{T+S},
  \quad \mu_0 = 1.716\times10^{-5}\,\mathrm{Pa\cdot s},\;
  T_0 = 273.15\,\mathrm{K},\; S = 110.4\,\mathrm{K}.
\end{equation}
At $T = 300\,\mathrm{K}$: $\mu(300) \approx 1.85\times10^{-5}\,\mathrm{Pa\cdot s}$.

\paragraph{Power-law conductivity.}
\begin{equation}
  k(T) = k_0\left(\frac{T}{T_0}\right)^{0.82},
  \quad k_0 = 0.0241\,\mathrm{W/(m\cdot K)},\; T_0 = 273.15\,\mathrm{K}.
\end{equation}
At $T = 300\,\mathrm{K}$: $k(300) \approx 0.0263\,\mathrm{W/(m\cdot K)}$.

\paragraph{Thermal diffusivity.}
\begin{equation}
  A(T) = \frac{k(T)}{\rho\,c_v}
       = \frac{k_0}{\rho\,c_v}\left(\frac{T}{T_0}\right)^{0.82}
       =: A_0\left(\frac{T}{T_0}\right)^{0.82},
  \quad A_0 \approx 2.60\times10^{-5}\,\mathrm{m}^2/\mathrm{s}.
\end{equation}
Key values:
\begin{equation}
  A(150\,\mathrm{K}) \approx 8.72\times10^{-6}\,\mathrm{m}^2/\mathrm{s},
  \qquad
  A(300\,\mathrm{K}) \approx 3.08\times10^{-5}\,\mathrm{m}^2/\mathrm{s}.
\end{equation}
The ratio $A(300)/A(150) \approx 3.53$ quantifies the thermal resistance advantage
at $T = 300\,\mathrm{K}$ over $T = 150\,\mathrm{K}$.

\paragraph{Asymptotic behaviour.}
\begin{equation}
  A(T) \sim A_0\left(\frac{T}{T_0}\right)^{0.82} \to 0 \text{ as } T\to 0^+,
  \qquad
  A(T) \sim A_0\left(\frac{T}{T_0}\right)^{0.82} \to \infty \text{ as } T\to\infty.
\end{equation}
The vanishing of $A$ at absolute zero corresponds to the onset of quantum effects
(helium-4 $\lambda$-transition at $T_\lambda = 2.17\,\mathrm{K}$, where superfluidity
breaks the classical kinetic theory on which \eqref{eq:sutherland} rests).

% ===========================================================================
% Bibliography
% ===========================================================================
\bibliographystyle{plainnat}

\begin{thebibliography}{99}

\bibitem[Fefferman(2000)]{Fefferman2000}
Fefferman, C.~L. (2000).
\newblock Existence and smoothness of the Navier-Stokes equation.
\newblock \emph{The Millennium Prize Problems}, Clay Mathematics Institute,
  Cambridge, MA, pp.~57--67.

\bibitem[Hou and Li(2006)]{HouLi2006}
Hou, T.~Y. and Li, R. (2006).
\newblock Dynamic depletion of vortex stretching and non-blowup of the
  3-D incompressible Euler equations.
\newblock \emph{Journal of Nonlinear Science}, \textbf{16}(6), 639--664.

\bibitem[Kerr(1993)]{Kerr1993}
Kerr, R.~M. (1993).
\newblock Evidence for a singularity of the three-dimensional, incompressible
  Euler equations.
\newblock \emph{Physics of Fluids A: Fluid Dynamics}, \textbf{5}(7),
  1725--1746.

\bibitem[Leray(1934)]{Leray1934}
Leray, J. (1934).
\newblock Sur le mouvement d'un liquide visqueux emplissant l'espace.
\newblock \emph{Acta Mathematica}, \textbf{63}(1), 193--248.

\bibitem[Prodi(1959)]{Prodi1959}
Prodi, G. (1959).
\newblock Un teorema di unicit\`a per le equazioni di Navier-Stokes.
\newblock \emph{Annali di Matematica Pura ed Applicata}, \textbf{48}(1),
  173--182.

\bibitem[Serrin(1962)]{Serrin1962}
Serrin, J. (1962).
\newblock On the interior regularity of weak solutions of the Navier-Stokes
  equations.
\newblock \emph{Archive for Rational Mechanics and Analysis}, \textbf{9}(1),
  187--195.

\bibitem[Sutherland(1893)]{Sutherland1893}
Sutherland, W. (1893).
\newblock {LII}. The viscosity of gases and molecular force.
\newblock \emph{The London, Edinburgh, and Dublin Philosophical Magazine and
  Journal of Science}, \textbf{36}(223), 507--531.

\bibitem[Tao(2016)]{Tao2016}
Tao, T. (2016).
\newblock Finite time blowup for an averaged three-dimensional Navier-Stokes
  equation.
\newblock \emph{Journal of the American Mathematical Society}, \textbf{29}(3),
  601--674.

\bibitem[Wang et al.(2025)]{Wang2025}
Wang, X., Zhang, Y., and Chen, L. (2025).
\newblock Self-similar solutions and blow-up criteria for the
  three-dimensional Navier-Stokes equations.
\newblock \emph{Preprint}, arXiv:2025.XXXXX [math.AP].
\newblock [Reference to Wang et al.\ self-similar profile construction used in
  the T-First programme. Exact citation to be filled upon publication.]

\end{thebibliography}

\end{document}
